\section{Compact convex envelopes and full-mass limits}
\label{sec:compact-convex-interfaces}
Compact effective domains give attainment for finite convex envelopes; preservation of total mass prevents measure loss at infinity.

\begin{lemma}\label{lem:compact-domain-finite-convex-envelope}
Let $G_1,\ldots,G_r$, with $r\geq1$, be closed proper convex functions whose effective domains lie in one compact set $K$. Their largest convex common minorant $H$ is closed and proper. At every point where $H$ is finite its finite-combination formula in \cref{prop:convex-envelope-algebra} attains its minimum. Consequently $(\max_iG_i^*)^*=H$. In particular, for finitely many $F_i\in\mathcal P(N_{\mathbb R},P)$, the convex envelope of their conjugates is already closed and equals $(\max_iF_i)^*$.
\end{lemma}
\begin{proof}
Each $G_i$ is bounded below on $K$. Otherwise a sequence with values tending to minus infinity would have a convergent subsequence there, contradicting lower semicontinuity and properness. Choose a common finite lower bound $B$. Convexity groups all terms involving $G_i$ into one term, so the envelope formula needs at most one point $x_i$ per function, with weights $\lambda_i\geq0$. Zero weights are omitted.

Take $x_\nu\to x$ with $\liminf H(x_\nu)=L<+\infty$, and choose approximate minimizing combinations along a subsequence realizing this limit. Their costs tend to $L$. Compactness of $K$ and the weight simplex gives a further subsequence with $x_{i,\nu}\to x_i$ and $\lambda_{i,\nu}\to\lambda_i$. For an absent term choose any fixed point of its effective domain. The limiting positive weights still sum to one and have barycenter $x$. The nonnegative quantities $\lambda_{i,\nu}(G_i(x_{i,\nu})-B)$ have bounded total. For $\lambda_i>0$, lower semicontinuity gives
\[
 \sum_{\lambda_i>0}\lambda_i(G_i(x_i)-B)
 \leq \liminf_\nu\sum_i
       \lambda_{i,\nu}(G_i(x_{i,\nu})-B)=L-B.
\]
In particular all those $G_i(x_i)$ are finite. Dropping the zero weights therefore gives an admissible combination of cost at most $L$. Thus $H(x)\leq L$, proving closedness. Also $H\geq B$ and $H\leq G_i$, so it is proper. Applying the same compactness argument to a minimizing sequence with fixed barycenter $x$ proves attainment whenever $H(x)<+\infty$. The conjugate identity and duality in \cref{prop:convex-envelope-algebra,thm:finite-dimensional-fenchel-duality} now give the formula. For the final assertion the conjugate domains lie in compact $P$ by \cref{prop:convex-support-growth-characterization}, and $F_i^{**}=F_i$.
\end{proof}
\index{convex envelope!compact effective domains}

\begin{lemma}\label{lem:full-mass-real-measure-bounded-testing}
Let finite positive Radon measures $\mu_j$ on a finite-dimensional real space converge vaguely to $\mu$. If $\mu_j$ and $\mu$ all have the same finite total mass, then $\int h\,\mathrm d\mu_j\to\int h\,\mathrm d\mu$ for every bounded continuous $h$. In particular this conclusion holds for the real Monge--Amp\`ere measures of locally uniformly convergent $F_j,F\in\mathcal E(N_{\mathbb R},P)$.
\end{lemma}
\begin{proof}
Choose $0\leq\chi\leq1$ continuous and compactly supported with $\int(1-\chi)\,\mathrm d\mu<\varepsilon$, using inner regularity and a cutoff equal to one on a compact set carrying all but $\varepsilon$ of the mass. Vague convergence of $\int\chi$ and equality of total masses give
\[
 \lim_j\int(1-\chi)\,\mathrm d\mu_j
       =\int(1-\chi)\,\mathrm d\mu<\varepsilon.
\]
The integrals of $\chi h$ converge. The two remaining tails have limiting sum at most $2\varepsilon\|h\|_\infty$, which proves the assertion by letting $\varepsilon\downarrow0$. For the real measures, \cref{thm:real-monge-ampere-vague-continuity} gives vague convergence and the common mass $n!\operatorname{vol}(P)$. When that mass is zero all measures vanish. Dimension zero also follows from this proof, with $\chi=1$ on the single point.
\end{proof}
\index{real Monge--Amp\`ere measure!full-mass convergence}
\index{tightness!from preservation of total mass}
